% !TEX root = ../Grunddatei.tex
\section{Material und Methoden}
Es wird aus einem Buch zitiert: Lorem ipsum\cite{Krom}

Beispiel f"ur Verwendung des Glossars: \Gls{latex} or \gls{latex}.

Beispiel f"ur Verwendung eines Symbols, dass im Symbolverzeichnis
auftaucht: \gls{Pi} oder \gls{N}.

Das erste Mal das ein Akronym wie \gls{FCN} benutzt wird, das zweite
mal \gls{FCN}, auch im Plural \glspl{FCN}. Man kann die Akronyme auf
wieder resetten...

\glsresetall

... dann geht es wieder von vorne los bei \gls{FCN}.